\documentclass[]{article}

\usepackage[T1]{fontenc}
\usepackage[utf8]{inputenc}
\usepackage{amsthm, amsmath, amsfonts}
\usepackage{blkarray}
\usepackage{bigstrut}
\usepackage{graphicx}

\newcommand{\avaliacao}{3}
\newcommand{\disciplina}{Probabilidade e Estatística}
\newcommand{\semestre}{2024/1}


\theoremstyle{definition}
\newtheorem{ex}{Exercício}

\renewcommand{\SS}{\mathcal{S}}

\begin{document}
\thispagestyle{empty}
\begin{center}
\Large{\avaliacao ª Avaliação de \disciplina\ - \semestre}

\large{Prof. Dr. Alessandro José Queiroz Sarnaglia}
\end{center}

\begin{ex}
Sejam $X_1, X_2, \ldots , X_{100}$ os pesos líquidos reais de 81 sacos de fertilizantes de 40 kg selecionados aleatoriamente.
\begin{enumerate}
	\item Se o peso esperado de cada saco for de fato 40 e a desvio-padrão 1,5, calcule $P(49\text{,}75 < \bar{X} < 50\text{,}25)$ (aproximadamente), usando o TLC;
	\item Se o peso esperado for 39,8 kg, e não 40 kg, de modo que em média os sacos tenham peso médio um pouco abaixo do anunciado um pouco a muito cheios, calcule $P(49\text{,}75 < \bar{X} < 50\text{,}25)$.
\end{enumerate}
\end{ex}


\begin{ex}
Seja $X_1, X_2, \ldots, X_n$ uma amostra aleatória de uma distribuição de Rayleigh com fdp
$$
f(x;\theta) = \frac{x}{\theta}\exp\left\{-\frac{x^2}{2\theta}\right\}, \ \ x > 0, \ \ \theta>0.
$$
\begin{enumerate}
	\item Encontre $E(X^2)$;
	\item Use o item acima para construir um estimador não viciado de $\theta$ com base em $\sum_iX_i^2$;
	\item Estime $\theta$ a partir das $n = 10$ observações a seguir sobre a força vibratória de uma palheta de turbina sob condições especificadas:
	\begin{center}
	\begin{tabular}{ccccc}
		16& 10& 14& 19&	4\\
		9 & 6 & 13& 6 &10\\
	\end{tabular}
	\end{center}
\end{enumerate}
\end{ex}


\begin{ex}
Em um experimento designado para medir o tempo adequado de exposição visual de um inspetor, a fim de realizar uma atividade com luz reduzida, o tempo médio da amostra para $n = 9$ inspetores foi 6,32 s e o desvio padrão da amostra 1,65 s. Assumiu-se anteriormente que o tempo de adaptação médio foi de pelo menos 7 s. Assumindo que o tempo de adaptação seja normalmente distribuído, os dados contradizem o ponto de vista anterior? Use o teste $t$ com $\alpha = 0,01$.
\end{ex}


\end{document}
